% !TeX root = ../../Book3.tex
\section{最大模}
\label{b3:sec:1901}

多项式与指数函数都在全平面全纯，却有不同的无穷远尺度。最大模把方向各异的增长压成一个径向量，三圆定理则限制这个量怎样随半径变化。

\subsection{\texorpdfstring{\(M(r,f)\)}{M(r,f)}}
\label{b3:subsec:190101}

对整函数 \(f\)，定义最大模 \(M(r,f)=\max_{|z|=r}|f(z)|\)。最大模原理使它等于闭圆盘上的最大模，并随 \(r\) 非减。若 \(f\) 非常值，则 \(M\) 严格递增：两半径相同最大值会在大圆盘内部取得最大模，迫使函数常数。

多项式 \(a_dz^d+\cdots+a_0\) 有 \(M(r)=|a_d|r^d(1+o(1))\)；\(e^z\) 有 \(M(r)=e^r\)。沿负实轴的指数趋零，与它的最大模指数增长并不矛盾。

\subsection{最大模}
\label{b3:subsec:190102}

\begin{proposition}\label{b3:prop:19-polynomial-growth}
若整函数满足 \(M(r,f)\le C(1+r)^N\)，则为次数不超过 \(\lfloor N\rfloor\) 的多项式。
\end{proposition}
\begin{proof}
Cauchy 估计给 \(|f^{(k)}(0)|\le k!C(1+r)^N/r^k\)。当整数 \(k>N\) 时令 \(r\to\infty\)，得到导数零；Taylor 展开因而有限。反向显然由有限项三角不等式成立。
\end{proof}
这里控制的是所有方向，而不是一条射线。函数 \(e^{-z}\) 在正实轴有界，却不是多项式。

\subsection{Hadamard 三圆定理}
\label{b3:subsec:190103}

\term[hadamard-sanyuan]{Hadamard 三圆定理}[Hadamard three-circles theorem]可由次调和比较直接得到。
\begin{theorem}[Hadamard 三圆]\label{b3:thm:19-three-circles}
若 \(f\) 在环域 \(r_1\le|z|\le r_3\) 邻域全纯且不恒零，则 \(r_1<r_2<r_3\) 时
\[
\log M(r_2)\le
\frac{\log(r_3/r_2)}{\log(r_3/r_1)}\log M(r_1)
+\frac{\log(r_2/r_1)}{\log(r_3/r_1)}\log M(r_3).
\]
\end{theorem}
\begin{proof}
选择 \(A,B\) 使调和函数 \(A\log|z|+B\) 在两边界等于相应最大对数。函数 \(\log|f|-A\log|z|-B\) 次调和，两边界均非正，最大值原理给内部非正。在半径 \(r_2\) 取最大值，线性插值就是题式。
\end{proof}

\subsection{对数凸性}
\label{b3:subsec:190104}

令 \(t=\log r\)，三圆定理说 \(t\mapsto\log M(e^t)\) 凸。其割线斜率随半径向外不减；这是增长约束，而不是通常变量 \(r\) 下的凸性。

例如单项式 \(z^d\) 给严格直线 \(dt\)；指数给 \(e^t\)。如果某非零整函数在所有大半径的对数斜率都不超过整数 \(d\)，则 \(M(r)\le Cr^d\)，上一命题将它识别为多项式。这使一个径向几何性质决定整个函数的代数结构。


\ProseParagraphBreak

三圆定理也能在同一半径上比较几何平均与最大值。Jensen 的圆周平均描述零点总影响，最大模则取最有利增长的方向；前者可能恒定，后者却增长很快。对 \(e^z\)，圆周对数平均零，而 \(\log M(r)=r\)。因此由平均对数的凸性不能直接推出最大模的精确增长，三圆定理另需对边界最大值作调和比较。这两种凸性虽然都以 \(\log r\) 为变量，控制的对象仍不同。

\ProseParagraphBreak

最大模关于半径连续。若 \(r_j\to r>0\)，在共同紧环带上 \(f\) 一致连续，将 \(r_je^{\mathrm{i}t}\) 与 \(re^{\mathrm{i}t}\) 比较，函数值差一致趋零，故两圆周最大值之差趋零。对非恒零整函数 \(M(r)>0\)，对数也连续。三圆不等式因此给一个通常意义的连续凸函数，而非只在排除零点的半径上定义的对象。圆周即使碰到函数零点，最大模仍由其他位置取得，不需要移开半径。

\ProseParagraphBreak

若三圆不等式在一个中间半径取得等号，比较函数 \(\log|f|-A\log|z|-B\) 在环域内部取得最大值零，故恒为零。于是 \(|f(z)|=e^B|z|^A\)，\(f\) 在环域无零点。局部对数微分给 \(f'/f=A/z\)，沿绕原点的圆周积分，辐角原理使 \(A\) 必为整数；继而 \(f(z)=cz^A\) 于环域。若 \(f\) 还在原点全纯，整数 \(A\) 必非负。三圆等号的刚性来自次调和强最大值原理，不能只由一个实凸函数的一次等号就辨认出复函数的形式。

\ProseParagraphBreak

最大模还可与 Taylor 展开的最大单项比较。对 \(f(z)=\sum c_nz^n\)，定义 \(\mu(r)=\max_n|c_n|r^n\)。整性使项随 \(n\to\infty\) 趋零，所以最大值确实达到。Cauchy 系数界给 \(\mu(r)\le M(r)\)；另一方面，对 \(0<r<R\)，
\[
 M(r)\le\sum_n|c_n|r^n
       \le\mu(R)\sum_n(r/R)^n
       =\frac{\mu(R)}{1-r/R}.
\]
这两边估计说明，模级数中的相消可能影响某个方向的函数值，却不能使整个圆周上的最大模比所有单项都小。取 \(R=2r\)，两种径向量具有相同的增长阶；当讨论有限正阶的型时，再取 \(R=(1+\varepsilon)r\) 并让 \(\varepsilon\downarrow0\)，它们的型也相同。
\ProseParagraphBreak
从对数坐标看，\(\log\mu(e^t)=\max_n(\log|c_n|+nt)\)，是斜率为非负整数的一族直线的上包络，因而凸。取得最大值的指标随 \(t\) 非减：若在 \(t_1<t_2\) 分别由 \(n_1,n_2\) 取到最大值，把两次比较不等式相加便得 \((n_2-n_1)(t_2-t_1)\ge0\)。这使 Taylor 展开的主导项随半径逐渐移向高次幂。若主导指标最终有共同上界，则所有足够高次系数都必须为零，否则该项的斜率更大，最终会超过已有的有限条直线。多项式与超越整函数的差别，由此再次表现为无穷远是否不断出现新的主导次数。
